\documentclass[11pt]{article}

\usepackage[a4paper, margin=2.4cm]{geometry}
\usepackage[T1]{fontenc}
\usepackage[utf8]{inputenc}
\usepackage{lmodern}
\usepackage{amsmath, amssymb, mathtools}
\usepackage{booktabs}

\title{LPV-SS vs.\ LPV-LFR Forward Pass}
\author{}
\date{}

\begin{document}
\maketitle

\noindent
Both methods compute the same trajectory. They differ in whether the
LFR structure $(G,\,\Delta(Y))$ is made explicit.

\bigskip

% -----------------------------------------------------------------------
% Side-by-side comparison
% -----------------------------------------------------------------------
\noindent
\begin{minipage}[t]{0.44\textwidth}
\centering\textbf{LPV-SS}

\medskip
\raggedright

\textbf{Step 1.} Net force:
\[ f_{\mathrm{net}} = -Kq - C\dot{q} + u \]

\textbf{Step 2.} Solve numerically at each time step:
\[ M(Y)\,a = f_{\mathrm{net}}, \quad a \in \mathbb{R}^3 \]

\textbf{Step 3.} Assemble state derivative:
\[ \dot{x} = \begin{bmatrix} \dot{q} \\ a \end{bmatrix} \]

\end{minipage}
\hfill
\vrule
\hfill
\begin{minipage}[t]{0.50\textwidth}
\centering\textbf{LPV-LFR}

\medskip
\raggedright

\textbf{Step 1.} Net force:
\[ f_{\mathrm{net}} = -Kq - C\dot{q} + u \]

\textbf{Step 2.} Set up the LFR loop equation:
\[ L(Y)\,z = \begin{bmatrix} M_0^{-1}f_{\mathrm{net}} \\ 0 \end{bmatrix},
   \quad L(Y) := I - D_{zw}\Delta(Y) \]

\textbf{Step 3a.} Solve block-by-block via Schur complement of $L(Y)$:$^*$
\[ z_2 = Y z_1, \qquad M(Y)\,z_1 = f_{\mathrm{net}} \]

\textbf{Step 3b.} Evaluate $M(Y)^{-1}$ as a rational function via Cramer's rule:$^{**}$
\[ z_1 = \frac{N(Y)\,f_{\mathrm{net}}}{d(Y)} =: a
   \quad\Rightarrow\quad
   z = \begin{bmatrix} a \\ Ya \end{bmatrix} \]

\textbf{Step 4.} Apply scheduling block:
\[ w = \Delta(Y)\,z = Yz \]

\textbf{Step 5.} State derivative through $G$:
\[ \dot{x} = A_x x + B_w w + B_u u \]

\end{minipage}

\medskip
\noindent\rule{\textwidth}{0.4pt}
\smallskip
\noindent
$^*$ Schur complement of the lower-right $I_3$ block in $L(Y)$ yields
$M_0^{-1}M(Y)$, giving $M(Y)z_1 = f_{\mathrm{net}}$ after cancelling $M_0^{-1}$.\\
$^{**}$ Cramer's rule applied to $M(Y)$ gives $M(Y)^{-1} = N(Y)/d(Y)$,
where $N(Y) = N_0 + YN_1 + Y^2N_2$ and $d(Y)$ are polynomial in $Y$,
precomputed once from physical parameters.
Both methods are used at different stages: Schur complement to solve the loop,
Cramer's rule to evaluate the result as a polynomial.

\bigskip

% -----------------------------------------------------------------------
% Properties table
% -----------------------------------------------------------------------
\begin{center}
\begin{tabular}{lcc}
\toprule
 & \textbf{LPV-SS} & \textbf{LPV-LFR} \\
\midrule
$G$ matrix explicit           & $\times$ & $\checkmark$ \\
$\Delta(Y)$ explicit          & $\times$ & $\checkmark$ \\
$z$, $w$ available            & $\times$ & $\checkmark$ \\
Runtime matrix solve          & $\checkmark$ ($3\times3$ per step) & $\times$ \\
Augmentation interface        & $\times$ & $\checkmark$ \\
Numerically equivalent        & \multicolumn{2}{c}{$\checkmark$} \\
\bottomrule
\end{tabular}
\end{center}

\end{document}
